%
% == INTRODUÇÃO ==
%
\chapter{Introdução}

O diagrama ER é uma representação gráfica usada para projetar e modelar um banco de dados. Ele auxilia na visualização das entidades (objetos, coisas ou conceitos) envolvidas em um sistema, bem como os relacionamentos entre essas entidades. 
Um diagrama ER é composto por entidades, atributos e relacionamentos. Segundo \cite{Fundamentos2017}, as entidades são representadas por retângulos e correspondem a objetos do mundo real, como pessoa, lugar, objeto e assim por diante. Os atributos são características dessa entidade, como nome, idade, entre outros. Os relacionamentos são representados por linhas que mostram como a entidade está conectada com outras.

Existem diferentes tipos de relacionamentos em um diagrama ER, como relacionamento um-para-um, um-para-muitos e muitos-para-muitos. Esses relacionamentos ajudam a definir como as entidades estão associadas entre si.
Além disso, é possível adicionar restrições e regras específicas ao diagrama, como chaves primária e estrangeira, que garantem a integridade dos registros no banco de dados.

Dada a rapidez com a qual os dados estão crescendo no cenário atual, a evolução da modelagem de um banco de dados tornou-se cada vez mais imperativa. Com o avanço da tecnologia e o aumento significativo na geração de dados, torna-se essencial que as ferramentas usadas para modelar esses bancos estejam alinhadas às demandas modernas.

No ambiente de desenvolvimento e projeto, a colaboração em tempo real tornou-se essencial, segundo \cite{Hsu2021}. Isso permite uma maior interação e comunicação entre os membros da equipe, independentemente da distância física. Com as ferramentas adequadas, é possível compartilhar informações de maneira eficaz, colaborar na criação de conteúdo e acompanhar o progresso do projeto em tempo real.

Essas ferramentas colaborativas agilizam a troca de ideias, o trabalho em equipe e a resolução de problemas \cite{Inoue2014}. Permitem, por exemplo, a realização de reuniões virtuais, o debate de ideias, o compartilhamento de documentos e atualizações instantâneas. Com isso, membros da equipe permanecem alinhados e informados, evitando atrasos e retrabalhos.

Ademais, a colaboração em tempo real promove a transparência e a visibilidade do projeto, uma vez que todos os envolvidos têm acesso à informação e podem acompanhar o progresso das atividades. Esse dinamismo contribui para um ambiente de trabalho mais eficiente, minimizando custos e otimizando a produtividade, conforme \cite{Fuks2002}.

 Portanto, o desenvolvimento de um modelador de diagramas ER colaborativo não é apenas oportuno, mas essencial para atender às demandas contemporâneas dos projetos de desenvolvimento de aplicações, garantindo que as soluções de banco de dados sejam robustas, eficientes e alinhadas com as necessidades de negócios em constante evolução.
 
%
% == JUSTIFICATIVA ==
%
\section{Justificativa}

Em um mundo digitalizado e interconectado, a excelência no projeto depende, em grande parte, da colaboração eficaz entre profissionais. A modelagem de um banco de dados, fundamental na arquitetura de um sistema, destaca-se nesse contexto. No entanto, o cenário atual revela uma carência de uma ferramenta acessível e colaborativa para a modelagem de diagramas ER. As alternativas disponíveis apresentam restrições, quer sejam de natureza funcional ou financeira. Por exemplo, algumas ferramentas, mesmo em seus pacotes mais elementares, possuem um custo elevado de aquisição e não disponibilizam todas as funcionalidades. Outras, em sua versão básica, não contemplam a opção para diagramas ER. Adicionalmente, a colaboração em tempo real, em certos casos, é limitada a um único usuário, relegando os demais a uma posição de mera visualização.

Diante da velocidade da inovação tecnológica e da complexidade crescente dos sistemas, uma ferramenta que permite a simultaneidade de contribuição de vários usuários é essencial. Ela não apenas potencializa o tempo de desenvolvimento, mas também favorece a identificação de inconsistências ou pontos de melhorias. Esta agilidade torna-se ainda mais crucial considerando a expansão da solução em nuvem e banco de dados distribuído, onde a colaboração remota é uma realidade inevitável.

O espectro dos beneficiários de tal ferramenta é vasto. Abrange desde um desenvolvedor ou arquiteto de sistema até o gestor de projeto, todos buscando um processo de modelagem mais ágil e coeso. Uma empresa com equipe dispersa geograficamente ou que opera sob metodologia ágil certamente verá um valor imensurável em uma plataforma que elimina barreira de comunicação e facilita a sinergia.

A longo prazo, a adoção de uma abordagem colaborativa na modelagem de diagramas ER tem o potencial de estabelecer um novo padrão na área. Com o reconhecimento de seu benefício, é provável que presenciemos uma transição para um método mais integrado, gerando transparência, eficácia e inovação. Ademais, à medida que uma tecnologia emergente, como a Inteligência Artificial, se integra ao processo, é plausível imaginar uma ferramenta que não apenas facilite a colaboração, mas também forneça sugestões e otimização automática, elevando a prática a um patamar inédito de sofisticação.

Considerando as demandas e desafios contemporâneos do mercado, o desenvolvimento de uma ferramenta colaborativa para modelagem de um diagrama ER não é apenas uma resposta adequada, mas uma iniciativa visionária, preparando o terreno para a necessidade futura da indústria.
	
%
% == OBJETIVOS ==
%
\section{Objetivos}

\subsection{Objetivo Geral}

Desenvolver um protótipo de modelador de diagramas ER colaborativo para a modelagem de bancos de dados, permitindo contribuições simultâneas de diversos usuários.

\subsection{Objetivos Específicos}	

\begin{enumerate}
    \item Modelar os dados de tabelas, relacionamentos, atributos, tipos de dados, posicionamento dos elementos na tela e outros. Além disso, persistir os dados em uma base de dados após a conclusão de uma sequência de ações.
    \item Investigar a biblioteca GoJS, compreendendo suas capacidades, limitações e possíveis aplicações no contexto do modelador de diagramas ER colaborativo proposto.
    \item Desenvolver uma interface que permita a múltiplos usuários trabalharem simultaneamente no mesmo diagrama ER, visualizando atualizações em tempo real.
    \item Desenvolver mecanismos de autenticação e autorização, garantindo que apenas usuários autorizados possam acessar e modificar os diagramas.
    \item Implementar práticas da metodologia ágil SCRUM e o método de gestão de fluxo Kanban durante todo o ciclo de desenvolvimento da aplicação.
\end{enumerate}
